\section{A process is more than a PID}
The process identifier is only a label. The meaningful idea is that each process has an execution context and a protected resource environment. Two instances of the same executable can therefore have different PIDs, stacks, heaps, open files, current directories, and scheduling states.

\subsection{Parent-child process relationships}
On Unix-like systems, processes form a hierarchy. A process created by \texttt{fork()} has a parent process. Immediately after a successful fork, parent and child begin with nearly identical logical memory contents, but they are separate processes. Modern systems normally implement this efficiently with \term{copy-on-write}: physical pages may initially be shared read-only and copied only when one process modifies a page.

\begin{keyidea}
\texttt{fork()} creates a new process; \texttt{exec()} replaces the current process image with a new program. Shells usually use them together: fork a child, let the child exec the requested command, and let the parent shell wait or continue depending on whether the job is foreground or background.
\end{keyidea}

\begin{mlfigure}
\begin{tikzpicture}[
  p/.style={draw=MLBlue,fill=MLPaleBlue,rounded corners=2mm,minimum width=3cm,minimum height=0.85cm,align=center,font=\small\bfseries\color{MLNavy}},
  arr/.style={-{Stealth[length=2mm]},draw=MLTeal,thick}
]
\node[p] (shell) {Shell process};
\node[p,below left=1.25cm and 1.2cm of shell] (parent) {Parent continues};
\node[p,below right=1.25cm and 1.2cm of shell] (child) {Child after fork};
\node[p,below=1.0cm of child] (prog) {New program after exec};
\draw[arr] (shell)--node[left,font=\scriptsize]{fork}(parent);
\draw[arr] (shell)--node[right,font=\scriptsize]{fork}(child);
\draw[arr] (child)--node[right,font=\scriptsize]{exec}(prog);
\end{tikzpicture}
\end{mlfigure}

\section{Waiting, zombies, and orphans}
When a child terminates, the kernel retains a small amount of exit information until its parent collects it using a wait operation. A terminated child whose status has not yet been collected is commonly called a \term{zombie}. A child whose parent exits first becomes an \term{orphan} and is adopted/re-parented according to the operating system's process-management rules.

\begin{pitfall}
A zombie is not a running process consuming CPU. It is a terminated process with an unreaped process-table entry. The resource leak is the retained bookkeeping entry, not continued execution.
\end{pitfall}

\section{File descriptors: per-process references to open objects}
A process usually has a table of small integer \term{file descriptors}. By convention, descriptors 0, 1, and 2 start as standard input, standard output, and standard error. An entry points to kernel-managed open-file state such as the current file offset and access mode.

\begin{workedexample}
If a shell redirects \texttt{program > out.txt}, the program itself can still write to descriptor 1 as usual. The shell has arranged for descriptor 1 to refer to \texttt{out.txt} before the program starts. This demonstrates why redirection is an OS/file-descriptor concept rather than a special feature inside every application.
\end{workedexample}

\section{Library calls versus system calls}
Not every C function that performs I/O maps one-to-one to a kernel transition. \texttt{fgetc()} is a C library function operating on a buffered \texttt{FILE *}. The library may satisfy multiple character reads from one larger underlying \texttt{read()} system call. Buffering reduces expensive transitions and improves throughput.

\section{Directory entries and pathname lookup}
A directory traversal with \texttt{readdir()} returns names and implementation-defined directory-entry information. The OS still has to resolve a pathname component by component. For a path such as
\[
\texttt{/home/student/report.txt},
\]
the file system starts from the root, looks up \texttt{home}, then \texttt{student}, then \texttt{report.txt}. Permissions and mount points can affect traversal.

\section{Defensive systems programming}
System interfaces can fail for ordinary reasons: a file does not exist, permissions are insufficient, memory is exhausted, a directory path is invalid, or a process limit is reached. Good systems code therefore treats error handling as part of the algorithm.
\begin{mlsteps}
  \item Check command-line argument counts before indexing \texttt{argv}.
  \item Check every resource-acquisition result such as \texttt{fopen()}, \texttt{opendir()}, \texttt{fork()}, and memory allocation.
  \item Use \texttt{perror()} or error codes to preserve the reason for failure.
  \item Release resources exactly once on every exit path.
  \item Test the return value of the read operation itself rather than predicting EOF in advance.
\end{mlsteps}

\begin{workedexample}
A robust text scan is conceptually:
\begin{enumerate}
  \item open the file;
  \item repeatedly request the next character;
  \item process it only if the read succeeded;
  \item stop when the read reports end-of-file or error;
  \item close the file.
\end{enumerate}
This is why \texttt{while ((ch = fgetc(fp)) != EOF)} is preferable to \texttt{while (!feof(fp))}.
\end{workedexample}

\begin{quickcheck}
\begin{enumerate}
  \item Why do shells commonly use both \texttt{fork()} and \texttt{exec()}?
  \item What is the difference between a \texttt{FILE *} and a POSIX integer file descriptor?
  \item Why can parent and child print in either order after \texttt{fork()}?
\end{enumerate}
\end{quickcheck}
